% ============================================================================
% 04-metodologia.tex — Metodologia (módulo 5/8)
% ============================================================================
\chapter{Metodologia}

\section{Desenho geral}

O estudo é observacional-experimental e auditável: quatro problemas canônicos
da IMO Shortlist (2020--2022) foram submetidos a LLMs gratuitas reais através
da CLI \texttt{opencode run}, em três condições experimentais, com rotação
controlada de contexto. O pipeline inteiro — da curadoria de paisagem ao
relatório — é versionado em Git (ciclos R497--R500) e validado por suíte de
testes (46 testes do R500 + 4066 anteriores, verdes na última rodada completa
pré-R500; regressão R500 verde em 1,97~s).

\subsection*{Curadoria acadêmica (R497)}
O \texttt{landscape/manifest.json} foi enriquecido com a fonte acadêmica
\texttt{deepmind/acme} (Apache-2.0), elevando o índice acadêmico para 5 fontes
(25 agentes indexados). A convergência semântica (\texttt{acme} $\leftrightarrow$
ecossistema) foi de 1,000 no laço de controle, 0,636 em agentes e pesquisa. A
terceira parte foi notificada em \texttt{THIRD\_PARTY\_NOTICES.md}.

\subsection*{Contrato de pipeline real (R498)}
O comando \texttt{/diagnose} foi regenerado com \texttt{deep=True}. O
\texttt{TestRealPipeContract} (3 testes) garante: (a) o template do comando
contém \texttt{deep=True}; (b) o módulo de integração idem; (c) execução real
de documentos. Execuções reais: ecossistema (13,8~s, 8 blocos),
\texttt{ARCHITECTURE.md} (SPS=1,0, NRR=0), \texttt{SPEC-935-R200}
(NRR=0,06, CR=0,8654), \texttt{07-discussion.tex} (CR=0,9712).

\section{Corpus IMO e solucionador determinístico}

\subsection*{Corpus}
% Problemas canônicos do harness (R442): algebra-001 (Shortlist 2021, resposta
% 3), algebra-004 (2^{u-2}), number-theory-001 (p³−q⁵=(p+q)², resposta (7,3),
% Shortlist 2022) e combinatorics-001 (2^{n−1}, n=1..6, Shortlist 2020).

Para a tarefa de benchmark usou-se o problema de teoria dos números
\texttt{imo-bench-number-theory-001}: encontrar primos $p,q$ tais que
$p^3 - q^5 = (p+q)^2$, cuja solução é $(p,q)=(7,3)$, verificação
$7^3 - 3^5 = 343 - 243 = 100 = 10^2 = (7+3)^2$.

\subsection*{Solucionador anti-leak (\texttt{RealIMOSolver}, R499)}
O \textit{harness} original (R442) ecoava a resposta esperada
(\texttt{short\_answer}) e atribuía nota 7 sozinho — \textit{leak} que
inviabilizava medição. O \texttt{RealIMOSolver} implementa solução
determinística: (i) enumeção de primos até 400 para o par $(7,3)$ (correto);
(ii) divisão inteira (\texttt{N=3}) versus racional (\texttt{N=1}) — achado
de ambiguidade de paráfrase; (iii) máximo local incluindo extremos
($2^{n-1}$) versus somente interior ($n=1,2$ falha); (iv) validação numérica
apenas, sem prova formal (status \textit{parcial}). Resultado: 2/4 corretos,
média 4,5/7 (alg001=6, nt001=6, alg004=3, comb001=3), além da sensibilidade
do \textit{grading head} à serialização ($2^{(n-1)}$ vs $2^{n-1}$).

\section{Condições experimentais}

Três condições (mesmo problema, mesma CLI, modelos free):
\begin{itemize}
  \item \textbf{C0 (crua)}: prompt direto sem contexto adicional;
  \item \textbf{C1 (contexto)}: prompt com contexto do problema (sem scaffold);
  \item \textbf{C2 (scaffold MASWOS)}: prompt com o template MASWOS de 16
  estágios, exigindo etapas explícitas (leitura, hipóteses, verificação,
  conformidade) e formato de saída estrito (um par).
\end{itemize}
A \textit{score} composto é $S = 0{,}40\,A + 0{,}30\,V + 0{,}20\,C + 0{,}10\,D$,
com acurácia $A$, velocidade normalizada $V = 1 - (\text{latência} - 20)/90$
(clampado em $[0,1]$), conformidade $C$ e disponibilidade $D$.

\section{Fluxograma do pipeline de benchmark}

% Fluxograma do processo (Figura 1): pasta fig-flux-pipeline.tex
% ============================================================================
% fig-flux-pipeline.tex — Fluxograma do pipeline de benchmark (Figura 1)
% ============================================================================
\begin{figure}[htbp]
\centering
\begin{tikzpicture}[
  node distance=0.9cm and 1.4cm,
  caixa/.style={rectangle, draw=cororo, thick, rounded corners=2pt,
                fill=cororo!8, minimum width=3.2cm, minimum height=0.9cm,
                align=center, font=\footnotesize},
  decisao/.style={diamond, draw=coramarelo, thick, fill=coramarelo!10,
                  aspect=2, align=center, font=\footnotesize},
  seta/.style={-{Stealth[length=3mm]}, thick, color=corcinza}
]
  \node[caixa] (prob) {Problema IMO\\ (nt001: $(7,3)$)};
  \node[caixa, right=of prob] (cond) {Condição\\ C0 / C1 / C2};
  \node[caixa, right=of cond] (llm) {LLM free real\\ \texttt{opencode run}};
  \node[decisao, below=of llm] (saida) {saída\\ conforme?};
  \node[caixa, below=of saida] (verif) {Verificação\\ \texttt{RealIMOSolver}};
  \node[caixa, below=of verif] (score) {Score composto\\ $S=0{,}4A+0{,}3V+0{,}2C+0{,}1D$};
  \node[caixa, below=of score] (rank) {Ranking\\ curadoria R499};

  \draw[seta] (prob) -- (cond);
  \draw[seta] (cond) -- (llm);
  \draw[seta] (llm) -- (saida);
  \draw[seta] (saida) -- node[right, font=\tiny]{sim} (verif);
  \draw[seta, color=corvermelho] (saida) -- ++(2.4,0) node[right, font=\tiny, color=corvermelho]{não: vazio/off-format $\rightarrow$ rejeita};
  \draw[seta] (verif) -- (score);
  \draw[seta] (score) -- (rank);
\end{tikzpicture}
\caption{Fluxograma do pipeline de \textit{benchmark} (R499--R500): o problema
IMO entra na condição C0/C1/C2; o LLM free responde via CLI; a saída é
verificada por conformidade e pelo solucionador determinístico; o \textit{score}
composto alimenta o ranking curado que dirige o roteamento R500.}
\label{fig:flux-pipeline}
\end{figure}

\section{Roteamento automático (R500)}

O \texttt{ModelRouter} foi estendido com catálogo free curado
(\texttt{integrations/free\_model\_catalog.py}) derivado do \textit{score}
R499: a rota \texttt{route\_free(task\_type)} seleciona o melhor modelo por
tipo de tarefa; \texttt{list\_all\_models()} passou a incluir
\texttt{opencode/big-pickle} e os cinco modelos free. O fluxograma do
roteamento é apresentado na seção de resultados (Figura 2).

\section{Replicação pós-R500}

Após a implementação (46 testes verdes), repetiu-se a condição C2 usando o
próprio roteamento (\texttt{route\_free("academic")} $\rightarrow$ mimo; n=3;
\texttt{route\_free("writing")} $\rightarrow$ big-pickle; n=2). A hipótese de
infraestrutura — estado acumulado do \textit{host} opencode de 1h15m sem
reinício — foi registrada e testada apenas parcialmente (reinício mataria a
sessão ativa); o não-determinismo de amostragem foi confirmado.

\section{Limitações pré-registradas}

(i) $n$ pequeno por condição; (ii) latência de infra variável (timeout
observado); (iii) ausência de reinício do \textit{host} antes das réplicas;
(iv) \textit{harness} de 4 problemas não é amostra estatística; (v) nenhuma
validação externa (editorial/benchmark independente) — portanto nenhum
resultado é rotulado \textit{superhuman} ou \textit{A1}.